% GENERATED FILE. Edit canonical lessons, then run npm run generate:books.
% canonical-source-hash: fnv1a64:02c62f74727b74b4
% canonical-lessons: GU-C37-kon, GU-C37-kyaan, GU-C37-kyaan-write, GU-C37-kyaare, GU-C37-ketla

\chapter{Who, Where, When, How Many}
\label{ch:gu-question-words}
\hlchaptermodality{41}

\begin{chapteropening}
\textbf{By the end of this chapter:} \emph{I can ask who, where, when and how many, and hear that every one of them is built on the same sound.}

The question family, and the pattern under it. Every Gujarati question word opens with ka-; the clause that answers opens with ja-; the thing pointed at opens with ta-. kyaare and jyaare differ by exactly one letter, so the answer taught two chapters ago turns out to have been half of a pair. kyaan finally lets the track ask where its own city, village, home, school and road are, and ketla lets it ask the price in a market it has taught since chapter 24.
\end{chapteropening}

\section[koṇ]{\gu{કોણ} (koṇ) — who}

\label{lesson:GU-C37-kon}



\begin{quote}
From the chapter behind you, the describing pair. And from chapter 3, the question word you have had longest.

\begin{itemize}
\raggedright
  \item \emph{Recall:} \emph{je … te …}, and \emph{kem}
\end{itemize}
\end{quote}

\begin{cousinweb}
\begin{quote}
\textbf{\gu{કોણ}} — \emph{koṇ} — \textbf{who}
\end{quote}

Look at the front of it. \textbf{\gu{કેમ}}, \textbf{\gu{કે}}, \textbf{\gu{કોણ}} — the same \textbf{\gu{ક}-} every time. Gujarati builds its question words on one stem, and once you have heard three of them the rest of the family announces itself.

This chapter is that family. Nothing in it is a new kind of thing; each lesson is one more question the same stem can ask.
\end{cousinweb}

\subsection*{Your turn}
\begin{itemize}
\raggedright
  \item \emph{Say it:} \emph{koṇ?} — flat, as a whole question.
  \item \emph{Contrast:} \emph{kem?} and \emph{koṇ?} — why, and who.
  \item \emph{Build:} ask who someone is.
\end{itemize}

\begin{tcolorbox}[breakable,colback=teal!4,colframe=teal!35!black,title={Before you move on}]
Ask \emph{who}. What do \textbf{\gu{કેમ}}, \textbf{\gu{કે}} and \textbf{\gu{કોણ}} share? (The \textbf{\gu{ક}-} at the front.)
\end{tcolorbox}
\section[kyā̃]{\gu{ક્યાં} (kyā̃) — where}

\label{lesson:GU-C37-kyaan}



\begin{quote}
One lesson back, \emph{who}. Then read the word for \textbf{city}, from the places chapter far behind you.

\begin{itemize}
\raggedright
  \item \emph{Read it:} \textbf{\gu{શહેર}}
\end{itemize}
\end{quote}

\begin{cousinweb}
\begin{quote}
\textbf{\gu{ક્યાં}} — \emph{kyā̃} — \textbf{where}
\end{quote}

The same \textbf{\gu{ક}-}, now with the \textbf{\gu{્ય}} you met inside \textbf{\gu{જ્યારે}}. The dot above is the anusvara: the vowel goes through the nose, and leaving it off makes a different word.

This is the question the track has most plainly been missing. It teaches \textbf{\gu{શહેર}}, \textbf{\gu{ગામ}}, \textbf{\gu{ઘર}}, \textbf{\gu{શાળા}} and \textbf{\gu{રસ્તો}} — city, village, home, school and road — and until this lesson there was no way to ask where any of them was.

Read that row aloud before you go on. Those five words were written more than a hundred lessons ago and most of them have not been asked for since; reading them cold, here, is the point of putting them in this lesson rather than a shorter example.
\end{cousinweb}

\subsection*{Your turn}
\begin{itemize}
\raggedright
  \item \emph{Say it:} \emph{kyā̃?} with the nasal audible.
  \item \emph{Build:} ask where a place you can name is.
  \item \emph{Contrast:} \emph{koṇ?} and \emph{kyā̃?} — a person, a place.
\end{itemize}

\begin{tcolorbox}[breakable,colback=teal!4,colframe=teal!35!black,title={Before you move on}]
Ask where the school is. What does the mark above \textbf{\gu{ક્યાં}} do? (Sends the vowel through the nose.)
\end{tcolorbox}
\section[kyā̃]{\gu{ક્યાં} reaches the page}

\label{lesson:GU-C37-kyaan-write}



\begin{quote}
One chapter back you wrote \textbf{\gu{તે}}. Write it once from memory.
\end{quote}

\subsection*{Script — four parts, none of them new}
Uncover \textbf{\gu{ક્યાં}} and take it apart in order:

\begin{quote}
\textbf{\gu{ક}} · the \textbf{\gu{્}} that binds it to · \textbf{\gu{ય}} · then the \textbf{\gu{ા}} stroke · and the \textbf{\gu{ં}} dot above
\end{quote}

Five signs, every one taught before chapter 30. This is the last word this run of chapters puts into the hand, and like every other one it spends nothing new.

\subsection*{Writing — guided copy, model in view}
Copy \textbf{\gu{ક્યાં}} once with the model in view. Check the dot is present — without it the word is not this word.

\begin{tcolorbox}[breakable,colback=teal!4,colframe=teal!35!black,title={Before you move on}]
Write \textbf{\gu{ક્યાં}} from memory and check the dot. How many of its signs were new today? (None.)
\end{tcolorbox}
\section[kyāre]{\gu{ક્યારે} (kyāre) — when}

\label{lesson:GU-C37-kyaare}



\begin{quote}
One lesson back you wrote \emph{where}. Two chapters back, the word that places one event against another.

\begin{itemize}
\raggedright
  \item \emph{Recall:} \textbf{\gu{ક્યાં}}, and \emph{jyāre}
\end{itemize}
\end{quote}

\begin{cousinweb}
Put them side by side. This is the neatest pattern in the chapter:

\begin{quote}
\textbf{\gu{ક્યારે}} — \emph{when?} · \textbf{\gu{જ્યારે}} — \emph{when …}
\end{quote}

\textbf{One letter apart.} The \textbf{\gu{ક}-} asks; the \textbf{\gu{જ}-} answers. You learned the answer two chapters ago without knowing it was half of a pair, and the question costs you almost nothing now.

That pairing runs through the whole language — \textbf{\gu{ક}-} for asking, \textbf{\gu{જ}-} for the clause that answers, \textbf{\gu{ત}-} for the thing pointed at. \textbf{\gu{જે} … \gu{તે} …} was the same system seen from the other side.
\end{cousinweb}

\subsection*{Your turn}
\begin{itemize}
\raggedright
  \item \emph{Contrast:} \emph{kyāre?} and \emph{jyāre} — question, then answer.
  \item \emph{Run through:} ask \emph{kyāre?}, then answer with a \emph{jyāre} clause.
  \item \emph{Say it:} the three fronts in a row — \emph{ka-}, \emph{ja-}, \emph{ta-}.
\end{itemize}

\begin{tcolorbox}[breakable,colback=teal!4,colframe=teal!35!black,title={Before you move on}]
Ask \emph{when}, then answer with a \emph{jyāre} clause. What single letter separates the two words? (The first — \textbf{\gu{ક}} asks, \textbf{\gu{જ}} answers.)
\end{tcolorbox}
\section[keṭlā]{\gu{કેટલા} (keṭlā) — how many}

\label{lesson:GU-C37-ketla}



\begin{quote}
The three questions this chapter has given you so far. Say them in order.

\begin{itemize}
\raggedright
  \item \emph{Recall:} \emph{koṇ}, \emph{kyā̃}, \emph{kyāre}
\end{itemize}
\end{quote}

\begin{cousinweb}
\begin{quote}
\textbf{\gu{કેટલા}} — \emph{keṭlā} — \textbf{how many}, \textbf{how much}
\end{quote}

The same \textbf{\gu{ક}-} one last time, now asking about quantity. This is the question the market chapter needed and never had: the track teaches \textbf{\gu{બજાર}}, \textbf{\gu{દુકાન}} and \textbf{\gu{પૈસા}} — market, shop and money — and could not ask the price of anything.

Four questions, one stem:

\begin{quote}
\textbf{\gu{કોણ}} who · \textbf{\gu{ક્યાં}} where · \textbf{\gu{ક્યારે}} when · \textbf{\gu{કેટલા}} how many
\end{quote}
\end{cousinweb}

\subsection*{Your turn}
\begin{itemize}
\raggedright
  \item \emph{Say it:} \emph{keṭlā?} — flat, as a whole question.
  \item \emph{Run through:} the four questions in a row, one breath each.
  \item \emph{Build:} ask how many of something you can name.
\end{itemize}

\begin{tcolorbox}[breakable,colback=teal!4,colframe=teal!35!black,title={Before you move on}]
Ask all four questions of this chapter. What do they share, and where did you first meet that shared piece? (The \textbf{\gu{ક}-}; in \textbf{\gu{કેમ}}, in chapter 3.)
\end{tcolorbox}
